\section{Exercise 3}

Các chuỗi nhắc nhập và thông báo được khai báo sẵn trong vùng \texttt{.data}:

\begin{lstlisting}[language=MIPS,numbers=none]
.data
  prompt_a: .asciiz "Insert a: "
  prompt_b: .asciiz "Insert b: "
  prompt_c: .asciiz "Insert c: "
  prompt_d: .asciiz "Insert d: "

  msg_f: .asciiz "F = "
  remainder_msg: .asciiz ", remainder "
\end{lstlisting}

Mỗi lần nhập in prompt bằng syscall 4, đọc số bằng syscall 5 rồi lưu vào thanh ghi. Dưới đây là khối nhập của $a$; ba biến còn lại làm tương tự với \texttt{prompt\_b}--\texttt{prompt\_d} và \texttt{\$s1}--\texttt{\$s3}:

\begin{lstlisting}[language=MIPS,numbers=none]
  li $v0, 4        # syscall 4: in prompt
  la $a0, prompt_a
  syscall
  li $v0, 5        # syscall 5: đọc số nguyên
  syscall
  move $s0, $v0    # lưu a vào $s0
\end{lstlisting}

Tử số gồm ba thừa số: $a+10$ dùng \texttt{addi}, $b-d$ dùng \texttt{sub}, còn $c-2a$ thì dịch trái $a$ 1 bit rồi trừ. Sau đó nhân dồn lại:

\begin{lstlisting}[language=MIPS,numbers=none]
  addi $t0, $s0, 10 # $t0 = a + 10
  sub $t1, $s1, $s3 # $t1 = b - d

  sll $t2, $s0, 1   # $t2 = 2a
  sub $t2, $s2, $t2 # $t2 = c - 2a

  add $t3, $s0, $s1 # $t3 = a + b
  add $t3, $t3, $s2 # $t3 = a + b + c

  mul $t0, $t0, $t1 # tử = (a + 10)(b - d)
  mul $t0, $t0, $t2 # nhân tiếp với (c - 2a)

  div $t0, $t3      # chia tử cho mẫu
  mflo $t0          # $t0 = thương
  mfhi $t1          # $t1 = số dư
\end{lstlisting}

Với $a=3$, $b=4$, $c=5$, $d=6$: tử số là $13 \times (-2) \times (-1) = 26$, mẫu số là $12$, nên $F = 2$ dư $2$.

\begin{lstlisting}[language={},numbers=none]
Insert a: 3
Insert b: 4
Insert c: 5
Insert d: 6
F = 2, remainder 2
\end{lstlisting}
